\chapter{Po bitke sú všetci\ldots}

Pôvodná predstava a zámer bakalárskej práce v počiatočnom štádiu výskumu a analýzy bol jednoduchý. Náš plán predstavoval
niekoľko krokov, ktoré tvorili základ problematiky:
\begin{enumerate}
  \item Identifikovať požiadavky na bezpečnosť ISVS vyplývajúce z legislatívy.
  \item Konkretizovať legislatívne požiadavky na základe noriem. 
  \item Navrhnúť postup, ktorým by sme efektívne naplnili tieto legislatívne požiadavky.
  \item Vytvoriť systém na podporu zavedenia ISMS\footnote{Zaviesť systém riadenia informačnej bezpečnosti od nás explicitne
    vyžadovala legislatíva.} v organizácii pre osobu poverenú touto úlohou -- manažéra KIB.
  \item Usporiadať a logicky zoskupiť informácie získane pri zavádzaní ISMS.
  \item Získané informácie využiť pri správe rizík.
\end{enumerate}

Počiatočný predpoklad bol taký, že na \enquote{malé} ISVS budú kladené iba základne (malé) bezpečnostné požiadavky a potrebný
systém na podporu by bol pomerne malý a jednoduchý. Ukázalo sa však, že opak bol pravdou a problematika sa v skutočnosti
ešte väčšmi rozšírila a vyšlo najavo, čo všetko sa v skutočnosti skrýva v požiadavkách zákonov. Nami navrhnutý systém je
síce nekompletný, nepokrýva všetky požiadavky stanovené zákonmi, no aj napriek tomu je veľmi rozsiahly.

Na svetlú stranu, bakalárska práca ukázala, že takýto systém na podporu vývoja bezpečnostného projektu sa dá vytvoriť:
\begin{enumerate}
  \item Podrobne sme analyzovali legislatívne požiadavky zákonov ZoKB a ZoITVS a vykonávacích predpisov 362/2018 a 179/2020
    v kapitole \ref{chapter:legal-analysis}.
  \item Zosúladili sme legislatívne požiadavky uvedených právnych noriem.
  \item Vyjadrili sme všeobecné požiadavky právnych noriem pomocou technických noriem ISO/IEC 27001 a ISO/IEC 27002\footnote{
      Z týchto dvoch technických noriem vychádzajú ZoKB a ZoITVS.}, s ktorými sú kompatibilné BSI 200-1,
    BSI 200-2, BSI 200-3\ a kompendium IT-Grundschutz.
  \item Upravili, resp. zjednodušili sme postup pri zavádzaní ISMS uvedený v metodikách BSI, aby bol zrozumiteľný aj pre
    neznalého a začínajúceho manažér KIB malého ISVS a najnižšej bezpečnostnej kategórie\footnote{Kategórie I z hľadiska
    kategorizácie sieti a informačných systémov (ZoKB) a správcu ISVS kategórie I (ZoITVS).}.
    KIB
  \item Vytvorili sme jednoduchý systém, ktorý prevedie manažéra KIB postupom zavádzania ISMS, poskytne mu základné informácie
    o KIB, legislatívnych požiadavkách a navrhovaných riešeniach.
\end{enumerate}
Nami vytvorený systém je dôkazom toho, že sa to dá, je to hľadaný \enquote{proof of concept} toho, že taký pomocný systém 
je možné vypracovať.

Nič však nie je dokonale a aj nášmu systému chýba mnoho ďalších vecí, na ktoré sme nemali dostatok času a boli mimo rámec 
bakalárskej práce. Preto by sme sa teraz radi venovali veciam, ktoré sme v písomnej práci alebo systéme nespomenuli, no
mali by byť súčasťou v nasledujúcich fáz\footnote{Poznamenáme, že zoznam určite nie je kompletný}, rozoberieme aj veci,
ktoré by sme mohli lepšie rozvrhnúť, lebo teraz už aj my sme generáli.

\section{Čo nebolo a mohlo byť, čo bolo a mohlo byť lepšie}

\subsection*{Ochrana osobných údajov}
Jednu z najväčších vecí, ktoré sme nespomenuli a sme sa jej vôbec nevenovali, je ochrana osobných údajov. Určite sme už všetci
počuli o nariadení 2016/679 Európskej únie, tiež známe ako GDPR\footnote{Po angl. \textit{General data protection regulation},
všeobecné nariadenie o ochrane údajov.} a vysokých sankciách udelených v prípade jeho porušenia. V organizácii by mala byť
prítomná špeciálna osoba, ktorá sa venuje ochrane osobných údajov. Určite by bolo potrebné vysvetliť túto problematiku v
systéme, no toto samo o sebe je rozsiahla oblasť hodná i svojej vlastnej záverečnej práce.

\subsection*{Príbuzná legislatíva}
V éteri sa tiež nachádza 241/2001 Z. z., zákon o ochrane utajovaných skutočností, náš systém by ho síce neriešil, ale jedna
časť by sa mu mohla venovať, aby uviedla manažéra do kontextu. V podobnom duchy by sme mohli informovať manažéra KIB o
zvyšných vykonávacích predpisoch ZoKB a ZoITVS, napríklad 493/2022 Z. z., o audite KIB alebo 78/2020 Z. z. o štandardoch
pre ITVS.

\subsection*{Bezpečnostné politiky druhej úrovne}
Politika KIB je všeobecný dokument, ktorý je síce rozsiahly, no nič špecifické nehovorí. Politiky druhej úrovne podrobnejšie
rozoberajú oblasti KIB obsiahnuté v politike KIB, ktoré sú rámcovo stanovené ZoKB a ZoITVS. Podrobne rozpracovať oblasti
je však veľmi rozsiahle.

\subsection*{Čo po spísaní politiky}
Spísať politiku KIB a zaviesť ISMS v organizácii je iba prvotný krok bezpečnostného procesu, v skutočnosti bude potrebné
udržiavať zavedený ISMS na dostatočnej úrovni, čím sa začne sa životný cyklus ISMS. Organizáciu ešte čaká vypracovanie bezpečnostných
projektov, inventarizácia a riadenie aktív, analýza a riadenie rizík.

\subsection*{Kategórie II a III}
Doposiaľ sme sa venovali iba najmenšej a z princípu najľahšej kategórií I. Síce podrobnú analýzu sme vykonali pre kategóriu
I, no nahliadli sme aj do vyšších kategórií, kde sú požiadavky na prevádzkovateľov základných služieb a správcov ISVS
rádovo väčšie. V krátkosti to znamená viac povinností, komplexnejšiu organizačnú štruktúru, väčší počet opatrení, zložitejšia
politika KIB a mnoho ďalšieho skrytého v legislatívnych požiadavkách.

\subsection*{Používatelia systému}
Pre náš systém je používateľ každý a súčasne nikto, chýba mu správa používateľov -- prihlasovanie sa. Prístup k systému
by potom mohli mať aj iní zamestnanci, pričom by sa nedostali k údajom, ktoré nie sú pre nich určené. Manažér KIB by mal
možnosť využiť systém na vzdelávanie alebo prípravu dotazníkov -- opýtaná osoba by sa prihlásila, dotazník vyplnila a
odpovede by boli priamo uložené a prístupné v systéme.

Je to aj otázka dostupnosti, čo ak by náš systém (alebo jeho derivácia) bola v budúcnosti úspešná? Malé ISVS by si nemuseli
  robiť starosti s inštaláciou systému, ba priam systém by bežal vo vl=dnom cloude. Na to ale treba vyriešiť autentifikáciu
a autorizáciu používateľov systému.

\subsection*{Upozornenia, správy a hlásenia}
Systém by mohol mať pomocnú funkcionalitu, ako hlásenie bezpečnostných incidentov pre zamestnancov\footnote{Koniec koncov
aj zákony vyžadujú zaviesť v organizácií jednotný systém hlásenia.}, evidencia bezpečnostných incidentov a generovanie
reportov za nejaké obdobie. Keby sa niekto pozabudol, systém by ho mohol v pravidelných intervaloch upozorniť, napr.
pripomenúť manažérovi KIB, že má neošetrené rizika.

\subsection*{Používateľská príručka}

\subsection*{Zamestnanci a role}

\subsection*{Kontrola ISMS, slovník a legislatíva}

\section{Zhrnutie}
